\documentclass[10pt,a4paper]{amsart}
\usepackage[margin=19mm]{geometry}
\usepackage{amsmath,amssymb,mathtools}
\usepackage[scr=rsfs,cal=euler]{mathalfa}
\usepackage{xfrac}
\usepackage[unicode,hidelinks,bookmarks=false]{hyperref}
\newcommand{\ie}{i.e.\ }
\newcommand{\cf}{cf.\ }
\newcommand{\resp}{respectively\ }
\newcommand{\Subsection}{\subsection}
\providecommand{\nobreakdash}{\nobreak\hbox{-}}
\setlength{\emergencystretch}{2em}
\multlinegap0pt
\usepackage{bm}
\usepackage{bbm}
\usepackage{dsfont}
\usepackage{xspace}
\usepackage{cancel}
\usepackage{mathtools}
\usepackage{amsthm}
\makeatletter
\@ifundefined{proposition}{\newtheorem{proposition}{Proposition}}{}
\makeatother
\newcommand\RR{\mathbb{R}}
\newcommand\ZZ{\mathbb{Z}}
\title[On existence of a compatible triangulation with the double c]{On existence of a compatible triangulation with the double circle order type}
\author{Hong Duc Bui}
\date{Source: arXiv:2508.04602v1 (CC BY 4.0)}
\begin{document}
\maketitle

\noindent\emph{Source and licence.} On existence of a compatible triangulation with the double circle order type. Version arXiv:2508.04602v1 (CC BY 4.0), distributed under the Creative Commons Attribution 4.0 International licence, as recorded in the arXiv licence metadata for this exact version and shown on the abstract page https://arxiv.org/abs/2508.04602v1. Full rights notice: licenses/arxiv-2508.04602-CC-BY-4.0.txt.\par\smallskip

\noindent\textit{Editorial scope.} Counted unit: the Proposition whose source label is \texttt{prop\_triangle\_subdivision\_3} in the article (source0.tex lines 464--476, source label \texttt{prop\_triangle\_subdivision\_3}), delivered together with the article's own complete original proof of it (source0.tex lines 477--569, one \texttt{proof} environment of 93 lines). No mathematics is written, repaired or re-derived: statement and proof are the article's own text line for line. Required context retained uncounted: none. Every definition, display and label of those lines is retained unchanged. Omitted from this package and not redistributed: 1--463, 570--2102. Nothing inside a retained excerpt is omitted or altered: every retained body is delivered exactly as the article's lines are written, so no omission receipt of any kind accompanies this package. Citations: the retained excerpts carry no citation command at all, so no bibliography entry of the article is transcribed and no reference is redistributed. Reference adaptation: none was needed and none was made; every reference inside the delivered unit resolves inside it, so no unresolved reference or citation is produced. Printed slips kept literal: no printed word, symbol, hypothesis or number of the counted statement or of its proof was altered. Layout and class: the article's own class is \texttt{article} with packages including amsmath, amssymb, biblatex, amsthm, xcolor, hyperref, cleveref, hypcap, graphicx, import, crossreftools; the wrapper is the lane's pinned \texttt{amsart} editorial wrapper at 10pt on A4, which restates the theorem-like environments the retained text needs and adds the article's own macro definitions that the retained text needs (\texttt{\textbackslash RR}, \texttt{\textbackslash ZZ}), with extra wrapper packages (bm, bbm, dsfont, xspace, cancel, mathtools). The delivered numbering is the unit's own. Assets: no retained span contains a figure environment, a graphics-inclusion command, a PDF-page inclusion, a table or a TikZ picture; no image file is copied into this package, the package declares no asset dependency and no image file is redistributed with it. Licence: version arXiv:2508.04602v1 of this article is distributed under the Creative Commons Attribution 4.0 International licence, as recorded in the arXiv licence metadata for this exact version and shown on the abstract page https://arxiv.org/abs/2508.04602v1. A full rights notice is delivered at licenses/arxiv-2508.04602-CC-BY-4.0.txt. Characters: every retained excerpt is pure ASCII, so the delivered engine, XeLaTeX with its default Latin Modern text font, reports no missing character.
\begin{proposition}
	\label{prop_triangle_subdivision_3}
	Given a convex polygon $P = P_1 P_2 \dots P_n$,
	points $\{A_1, \dots, A_m \}$ inside polygon $P$,
	such that $P \cup \{ A_1, \dots, A_m \}$ is in general$^+$ position.
	Let $c_t$ and $c_b$ be nonnegative integers, $c_i$ be a positive integer,
	such that $c_t + c_i + c_b = m$.
	Then there exists a point $Q$ inside triangle $P$ such that within the $m$ points
	$\{ A_1, \dots, A_m \}$,
	there are exactly $c_t$ points inside triangle $P_1 P_i Q$,
	$c_b$ points inside triangle $P_1 Q P_{i+1}$,
	and $c_i$ points inside triangle $Q P_i P_{i+1}$.
\end{proposition}

\begin{proof}
	\newcommand\bracket[1]{\langle #1 \rangle}
	The idea is the following.
	Consider the function $\hat f \colon P \to \ZZ^3$ that maps any point inside triangle $P$ (including the boundary)
	to a triple of integers, defined by $\hat f(Q) = (\bracket{P_1 P_i Q}, \bracket{P_1 Q P_{i+1}}, \bracket{Q P_i P_{i+1}})$,
	where for a triangle $T$, $\bracket{T}$ is the number of points within $\{ A_1, \dots, A_m \}$ inside that triangle.

	Then we wish to find some point $Q$ such that $\hat f(Q) =(c_t, c_b, c_i)$.
	To do that, we find a continuous function $f$ that approximates $\hat f$,
	then apply a result from algebraic topology.

	\renewcommand\bracket[1]{[#1]}

	Let $\varepsilon>0$ be some very small positive integer, which will be selected later.
	For each point $A_j$, let $B(A_j, \varepsilon)$ be the closed ball with radius $\varepsilon$ centered at $A_j$.
	Let $V = \varepsilon^2 \pi$ be the area of such a ball.

	For a geometric object $D$ on the plane, define $|D|$ to be the area of $D$.

	Define a function $f \colon P \to \RR^3$ as follows.
	For each point $Q \in P$,
	\[ f(Q) = (\bracket{P_1 P_i Q}, \bracket{P_1 Q P_{i+1}}, \bracket{Q P_i P_{i+1}}), \]
	where for a triangle $T$, define \[ \bracket{T} = \frac{1}{m \cdot V} \sum_{j=1}^m |B(A_i, \varepsilon) \cap T|. \]

	It is clear that as $\varepsilon \to 0$, $f$ converges to $\hat f$ almost everywhere.
	This is what we mean by $f$ being an approximation of $\hat f$.

	Assume $\varepsilon$ is small enough such that for each $1 \leq j \leq m$,
	$B(A_j, \varepsilon) \subseteq P$.
	Then for each point $Q \in P$, the sum of coordinates of $f(Q)$ is exactly $m$.

	Define $\Delta^2$ to be the closed affine triangle in $\RR^3$ with
	three vertices having coordinate $(1, 0, 0)$, $(0, 1, 0)$ and $(0, 0, 1)$.
	Using the argument above, combined with the fact that each coordinate of $f$ is nonnegative,
	we have the image of $f$ is inside $\Delta^2$.
	We thus restrict the codomain of $f$ to $\Delta^2$.

	Furthermore, we have the following:
	\begin{itemize}
		\item $f(P_1)=(0, 0, m)$,
		\item for each point $Q$ on segment $P_1 P_i$, $f(Q)$ has first coordinate $0$,
		\item $f(P_i)=(0, m, 0)$,
		\item for each point $Q$ on segment $P_i P_{i+1}$, $f(Q)$ has third coordinate $0$,
		\item $f(P_{i+1})=(m, 0, 0)$,
		\item for each point $Q$ on segment $P_{i+1} P_1$, $f(Q)$ has second coordinate $0$.
	\end{itemize}
	Therefore, as $Q$ travels once around the boundary of $P$,
	$f(Q)$ travels once around the boundary of $\Delta^2$.
	In the language of algebraic topology, $f$ maps the boundary of $P$ to the boundary of $\Delta^2$,
	and the induced maps on the first homology group $H_1$ and the fundamental group $\pi_1$ are bijective.

	Therefore, $f \colon P \to \Delta^2$ is surjective,
	using an argument similar to that used in the proof of Brouwer fixed-point theorem.
	This gives us some point $Q$ such that $f(Q) = (c_t, c_b, c_i)$.

	Recall that we want to find $Q$ such that $\hat f(Q) =(c_t, c_b, c_i)$, so we are almost there.

	It remains to prove that for sufficiently small $\varepsilon$ there exists a point $Q$ such that $f(Q) = (c_t, c_b, c_i)$,
	and also $\hat f(Q) = (c_t, c_b, c_i)$.
	We see that it suffices if none of the segments $P_1 Q$, $P_i Q$ and $P_{i+1} Q$
	has any intersection with any of the balls $B(A_j, \varepsilon)$---%
	informally, none of the balls are cut by the segments that separate the three triangles.

	First, we rule out the possibility that $Q$ is inside some ball $B(A_j, \varepsilon)$.
	If this is the case, because no three points are collinear,
	for sufficiently small $\varepsilon$,
	none of the segments $P_1 Q$, $P_i Q$, or $P_{i+1} Q$
	can intersect with any other ball $B(A_k, \varepsilon)$ for $k \neq j$.
	This makes $f$ have non-integral coordinate because $B(A_j, \varepsilon)$ is split into multiple parts,
	so $f(Q) \neq (c_t, c_b, c_i)$, contradiction.

	Then we rule out the possibility that all three segments
	$P_1 Q$, $P_i Q$, and $P_{i+1} Q$
	intersect some balls
	$B(A_{j_1}, \varepsilon)$,
	$B(A_{j_2}, \varepsilon)$,
	and $B(A_{j_3}, \varepsilon)$ simultaneously.
	If all three are the same ball, that is $j_1 = j_2 = j_3$, then point $Q$ must be inside the ball,
	which is already ruled out.
	If two of them are the same ball, let's say $j_1 = j_2$,
	using the same argument as above, because no three points are collinear,
	for sufficiently small $\varepsilon$,
	there cannot be any other $B(A_{j_3}, \varepsilon)$ lying on the remaining segment.
	The remaining case is if all $j_1$, $j_2$ and $j_3$ are distinct,
	then because no three lines are concurrent, for sufficiently small $\varepsilon$ this does not happen either.

	The last case is if there are at most two segments
	$P_1 Q$, $P_i Q$, and $P_{i+1} Q$
	intersecting some ball.
	Again, because no three points are collinear,
	each segment intersects at most one ball for small enough $\varepsilon$.
	Therefore, $f(Q)$ has some nonintegral coordinate, contradiction.
\end{proof}

\end{document}
